\documentclass[10pt,a4paper]{amsart}
\usepackage[margin=19mm]{geometry}
\usepackage{amsmath,amssymb,mathtools}
\usepackage[scr=rsfs,cal=euler]{mathalfa}
\usepackage{xfrac}
\usepackage[unicode,hidelinks,bookmarks=false]{hyperref}
\newcommand{\ie}{i.e.\ }
\newcommand{\cf}{cf.\ }
\newcommand{\resp}{respectively\ }
\newcommand{\Subsection}{\subsection}
\providecommand{\nobreakdash}{\nobreak\hbox{-}}
\setlength{\emergencystretch}{2em}
\multlinegap0pt
\usepackage{bm}
\usepackage{bbm}
\usepackage{dsfont}
\usepackage{xspace}
\usepackage{cancel}
\usepackage{mathtools}
\usepackage{amsthm}
\makeatletter
\@ifundefined{thm}{\newtheorem{thm}{Thm}}{}
\@ifundefined{lem}{\newtheorem{lem}[thm]{Lemma}}{}
\@ifundefined{prop}{\newtheorem{prop}[thm]{Proposition}}{}
\makeatother
\newcommand*{\EE}{\ensuremath{\mathbb{E}}}	        	% Erwartung
\newcommand*{\HH}{\ensuremath{\mathbb{H}}}	        	% harmonic functions
\newcommand*{\II}{\ensuremath{\mathds{1}}}	        	% Indikatorfunktion
\newcommand*{\PP}{\ensuremath{\mathbb{P}}}	        	% W-Maß
\newcommand*{\RR}{\ensuremath{\mathbb{R}}}	        	% reelle Zahlen
\newcommand*{\eps}{\ensuremath{\varepsilon}}            % das bessere epsilon
\title[Hyperbolic branching Brownian motion: the empirical limit me]{Hyperbolic branching Brownian motion: the empirical limit measure}
\author{David Geldbach}
\date{Source: arXiv:2509.06730v3 (CC BY 4.0)}
\begin{document}
\maketitle

\noindent\emph{Source and licence.} Hyperbolic branching Brownian motion: the empirical limit measure. Version arXiv:2509.06730v3 (CC BY 4.0), distributed under the Creative Commons Attribution 4.0 International licence, as recorded in the arXiv licence metadata for this exact version and shown on the abstract page https://arxiv.org/abs/2509.06730v3. Full rights notice: licenses/arxiv-2509.06730-CC-BY-4.0.txt.\par\smallskip

\noindent\textit{Editorial scope.} Counted unit: the Prop whose source label is \texttt{prop:2nd\_moment} in the article (source0.tex lines 645--661, source label \texttt{prop:2nd\_moment}), delivered together with the article's own complete original proof of it (source0.tex lines 665--726, one \texttt{proof} environment of 62 lines). No mathematics is written, repaired or re-derived: statement and proof are the article's own text line for line. The proof is detached from the statement in the source: the article states the result at lines 645--661 and prints its proof at lines 665--726, and the proof environment's own header names the statement, so the association is the article's own. Required context retained uncounted: lemma:many\_to\_two (lines 602--612); lem:bound\_on\_bm (lines 624--630). Every definition, display and label of those lines is retained unchanged. Omitted from this package and not redistributed: 1--601, 613--623, 631--644, 662--664, 727--928. Nothing inside a retained excerpt is omitted or altered: every retained body is delivered exactly as the article's lines are written, so no omission receipt of any kind accompanies this package. Citations: the retained excerpts carry no citation command at all, so no bibliography entry of the article is transcribed and no reference is redistributed. Reference adaptation: none was needed and none was made; every reference inside the delivered unit resolves inside it, so no unresolved reference or citation is produced. Printed slips kept literal: no printed word, symbol, hypothesis or number of the counted statement or of its proof was altered. Layout and class: the article's own class is \texttt{article} with packages including inputenc, xcolor, fontenc, babel, graphicx, latexsym, amsmath, amssymb, amsthm, commath, enumitem, mathrsfs, geometry, dsfont; the wrapper is the lane's pinned \texttt{amsart} editorial wrapper at 10pt on A4, which restates the theorem-like environments the retained text needs and adds the article's own macro definitions that the retained text needs (\texttt{\textbackslash EE}, \texttt{\textbackslash HH}, \texttt{\textbackslash II}, \texttt{\textbackslash PP}, \texttt{\textbackslash RR}, \texttt{\textbackslash eps}), with extra wrapper packages (bm, bbm, dsfont, xspace, cancel, mathtools). The delivered numbering is the unit's own. Assets: no retained span contains a figure environment, a graphics-inclusion command, a PDF-page inclusion, a table or a TikZ picture; no image file is copied into this package, the package declares no asset dependency and no image file is redistributed with it. Licence: version arXiv:2509.06730v3 of this article is distributed under the Creative Commons Attribution 4.0 International licence, as recorded in the arXiv licence metadata for this exact version and shown on the abstract page https://arxiv.org/abs/2509.06730v3. A full rights notice is delivered at licenses/arxiv-2509.06730-CC-BY-4.0.txt. Characters: every retained excerpt is pure ASCII, so the delivered engine, XeLaTeX with its default Latin Modern text font, reports no missing character.
\begin{lem}[many-to-two] \label{lemma:many_to_two}
    {Fix $r\geq 0$,} under $\PP^r$, let $(X_s^1,Y_s^1)_{s\geq 0}$ and $(X_s^2,Y_s^2)_{s\geq 0}$ be two hyperbolic Brownian motions that move together until time $r$ and afterwards move independently. Then we have for any interval $I\subseteq \RR$ that 
    \begin{equation*}
        \EE\left[\mu_\infty(I)^2\right] = 2\beta\int_0^\infty\PP^r\left(X^1_\infty,X_\infty^2\in I\right) e^{-\beta r}dr.
    \end{equation*}
    Similarly for any $K>0$,
        \begin{align*}
        &\EE\left[\mu_\infty^K(I)^2\right] = 2\beta\int_0^\infty\PP^r\bigg(X^1_\infty,X_\infty^2\in I, \\ 
        & \hspace{4cm}\forall s\geq K: \log(Y_s^1)+s/2, \log(Y_s^2)+s/2 \in [-s^{2/3},s^{2/3}]\bigg) e^{-\beta r}dr.
    \end{align*}
\end{lem}

\begin{lem}\label{lem:bound_on_bm}
    For any $x\in \RR$, $y\in (0,\infty)$ and any interval $I\subseteq \RR$ we have that
    \begin{equation*}
        \PP_{(x,y)} (X_\infty\in I)\leq\left( \frac{1}{\pi} \frac{\vert I \vert}{y} \right)\wedge 1,
    \end{equation*}
    {where $\vert I \vert$ is the length of the interval.}
\end{lem}

\begin{prop}\label{prop:2nd_moment} $\hspace{1cm} $
    \begin{enumerate}
        \item [(i)]
    For any $\beta>0$ and $\beta \neq 3$ there is $C_1<\infty$ such that
    \begin{equation*}
        \limsup_{\eps \to 0} \sup_{\substack{I\subseteq \RR \\ \vert I \vert =\eps}} \frac{\EE\left[ \mu_\infty(I)^2\right]}{\eps^{2 \wedge(1+\beta/3)}} \leq C_1<\infty.
    \end{equation*}
    For $\beta=3$, replace $\eps^2$ above by $\eps^2 \log(\eps^{-1})$.

    \item [(ii)]
    For any $\beta>0 $, any $K>0$, and any $\delta>0$, there is $C_2=C_2(\beta, K, \delta)<\infty$ such that 
    \begin{equation*}
        \limsup_{\eps \to 0} \sup_{\substack{I\subseteq \RR \\ \vert I \vert =\eps}} \frac{\EE\left[ \mu^K_\infty(I)^2\right]}{\eps^{2 \wedge(1+2\beta-\delta)}} \leq C_2<\infty.
    \end{equation*}

    \end{enumerate}
\end{prop}

\begin{proof}[Proof of Proposition \ref{prop:2nd_moment} (i).] Assume for now that $\beta\neq 3$.
    Without loss of generality, we consider $I=[-\eps/2,\eps/2]$. 
    
    We use Lemma \ref{lemma:many_to_two} and subdivide the integral into three parts, $J_1=[0,1],J_2=[1,\log(\eps^{-2})]$ and $J_3=[\log(\eps^{-2}), \infty)$. For $i\in \{1,2,3\}$, let
\begin{equation*}
    T_i = \int_{J_i} \PP^r\left(X^1_\infty,X_\infty^2\in I\right) e^{-\beta r}dr.
\end{equation*}
We start with $T_1$. Here we bound $e^{-\beta r} \leq 1$ and then condition on the splitting position at time $r$,
\begin{align*}
    T_1 = \int_0^1 \int_{\HH}  \PP_{(x,y)}\left(X_\infty \in I\right)^2 \PP\left((X_r,Y_r) \in (dx,dy)\right) dr.
\end{align*}
We apply Lemma \ref{lem:bound_on_bm} {to obtain}
\begin{equation}\label{eq:T1}
    T_1 \leq c_1 \eps^2\int_0^1 \EE\left[Y_r^{-2}\right]dr \leq c_2 \eps^2,
\end{equation}
for some {constants} $c_1,c_2>0$. 

Next we consider $T_2$. By Lemma \ref{lem:bound_on_bm},
\begin{align*}
    \PP^r\left(X^1_\infty,X_\infty^2\in I\right) &=\int_\HH \PP((X_r,Y_r)\in(dx,dy))\PP_{(x,y)}(X_\infty \in I)^2 \\
    &\leq \int_\HH \PP((X_r,Y_r)\in(dx,dy))\PP_{(x,y)}(X_\infty \in I)\PP_{(r,0,y)}(X_\infty \in I)\\
    &\leq \int_\HH \PP((X_r,Y_r)\in(dx,dy))\PP_{(x,y)}(X_\infty \in I)\left(\frac{1}{\pi}\frac{\eps}{y} \wedge 1\right).
\end{align*}
{Hence}
\begin{align*}
    T_2 \leq \int_1^{\log(\eps^{-2})}\EE\left[\II_{X_\infty \in I} \left(\frac{1}{\pi}\frac{\eps}{Y_r} \wedge 1\right)\right] e^{-\beta r}dr.
\end{align*}
We split this integral into two parts: for $r\in [1,\log(\eps^{-1/3})]$ we bound
\begin{align*}
    \EE\left[ \II_{X_\infty\in I}\left(\frac{1}{\pi}\frac{\eps}{Y_r} \wedge 1\right)\right] &\leq \frac{1}{\pi}\eps\EE\left[\II_{X_\infty\in I}  Y_r^{-1}\right] = \frac{1}{\pi}\eps \EE\left[\PP(X_\infty\in I \vert Y_r) 
    Y_r^{-1}\right]\\
    &\leq \frac{1}{\pi^2}\eps^2 \EE\left[
    Y_r^{-2}\right] = \frac{1}{\pi^2}\eps^2 e^{3r},
\end{align*}
where we used Lemma \ref{lem:bound_on_bm} again, and where we recall that $Y_r^{-1} = \exp(r/2-B_r)$.
{F}or $r\in  [\log(\eps^{-1/3}),\log(\eps^{-2})]$, {we estimate}
\begin{equation*}
    \EE\left[ \II_{X_\infty\in I}\left(\frac{1}{\pi}\frac{\eps}{Y_r} \wedge 1\right)\right] \leq \PP(X_\infty \in I) \leq \frac{1}{\pi}\eps,
\end{equation*}
where we also used Lemma \ref{lem:bound_on_bm}.
Then our estimate on $T_2$ becomes
\begin{align}
    T_2 &\leq \frac{1}{\pi^2}\eps^2\int_1^{\log(\eps^{-1/3})} e^{3r}e^{-\beta r}dr + \frac{1}{\pi}\eps\int_{\log(\eps^{-1/3})}^{\log(\eps^{-2})}e^{-\beta r}dr \nonumber\\
    &\leq c_3\left(\eps^2 + \eps^{1+\beta/3} + \eps^{1+2\beta}\right), \label{eq:T2}
\end{align}
for some $c_3>0$.
For $T_3$, we estimate using Lemma \ref{lem:bound_on_bm}
\begin{equation*}
    \PP^r(X_\infty^1, X_\infty^2 \in I)\leq \PP(X_\infty^1\in I) \leq \frac{\eps}{\pi}.
\end{equation*}
Therefore
\begin{align}\label{eq:T3}
    T_3\leq \frac{1}{\pi}\eps \int_{\log(\eps^{-2})}^\infty e^{-\beta r}dr = c_4 \eps^{1+2\beta},
\end{align}
for $c_4>0$. 
To complete the proof, we combine \eqref{eq:T1}, \eqref{eq:T2} and \eqref{eq:T3}.

Lastly, in the case where $\beta=3$, the final estimate on $T_2$ is 
\begin{align}
    T_2 &\leq \frac{1}{\pi^2}\eps^2\int_1^{\log(\eps^{-1/3})} e^{3r}e^{-3 r}dr + \frac{1}{\pi}\eps\int_{\log(\eps^{-1/3})}^{\log(\eps^{-2})}e^{-3 r}dr \nonumber\leq c_3 \eps^2 \log(\eps^{-1}). 
\end{align}
\end{proof}

\end{document}
